% 第 4 讲：整体类域论与二次互反律
% 源：src/lecture04.txt（379 行，PDF 抽出文本，公式按数学含义重排）
% 块：R2 / reports/R-block02.md

\section*{第 4 讲\quad 整体类域论与二次互反律}

上一讲中，我们把希尔伯特符号 $(\cdot,\cdot)$ 的非退化性与双乘性，归结为证明：对一切二次扩张 $L/K$（$K$ 为局部域），$NL^{\times}\subseteq K^{\times}$ 都是指数为 $2$ 的子群。我们已经证明，当扩张非分歧、以及当 $p=\operatorname{char}(\OK_K/\mathfrak{p})$ 为奇数时结论成立（$p=2$ 的情形已在习题集 1 第 1 题中大致证明）。本讲中，我们将在整体情形，即 $F=\Q$ 的情形，做类似的分析。

我们比较以下两个事实。第一，
\[
\Gal^{\mathrm{ab}}(\Q)/2\cong \operatorname{Hom}\bigl(\Q^{\times}/(\Q^{\times})^2,\{1,-1\}\bigr)=\operatorname{Hom}(\Q^{\times},\{1,-1\}),
\]
其中全体素数连同 $-1$ 构成这个群作为 $\F_2$-向量空间的一组基。第二，cft（类域论）预测
\[
\mathbb{A}^{\times}_{\Q}/\Q^{\times}\cong \Gal^{\mathrm{ab}}(\Q)
\]
是投射有限完备化之间的典范同构。于是我们期待
\begin{equation*}\tag{4.1}
\operatorname{Hom}(\Q^{\times},\{1,-1\})=\Gal^{\mathrm{ab}}(\Q)/2\cong \Q^{\times}\backslash \mathbb{A}^{\times}_{\Q}/(\mathbb{A}^{\times}_{\Q})^2 .
\end{equation*}

回忆下面的定义：

\textbf{定义 4.1.} \emph{阿代尔环}定义为
\[
\mathbb{A}_{\Q}:=\varinjlim_{S}\Bigl(\prod_{p\notin S}\Z_p\times\prod_{p\in S}\Q_p\Bigr),
\]
其中 $S$ 取遍 $\Q$ 的有限素位集合（素数与 $\infty$），按包含关系定向，并约定 $\Q_{\infty}=\R$。\emph{伊代尔群}是阿代尔环的乘法群
\[
\mathbb{A}^{\times}_{\Q}=\varinjlim_{S}\Bigl(\prod_{p\notin S}\Z^{\times}_p\times\prod_{p\in S}\Q^{\times}_p\Bigr),
\]
其中 $S$ 同上。
% [原文疑误] 抽出的伊代尔群定义中第一个乘积印作 ∏_{p∈S} Z_p^×，与阿代尔环的定义不符，应为 p∉S；已按后者排。

\textbf{例 4.2.} 我们有 $\widehat{\Z}=\prod_p \Z_p$，其中 $p$ 遍历所有素数；而 $\widehat{\Z}\times\R\subseteq \mathbb{A}_{\Q}$ 作为子环嵌入，其中可以把任意 $n\neq 0$ 对角嵌入。类似地，$\widehat{\Z}^{\times}=\prod_p \Z^{\times}_p$，且 $\widehat{\Z}^{\times}\times\R^{\times}\subseteq \mathbb{A}^{\times}_{\Q}$。然而 $2$ 与 $1/2$ 都不在 $\widehat{\Z}^{\times}\times\R^{\times}$ 中，因为它们不在 $\Z^{\times}_2$ 中；任何有理数都是如此。若我们把有理数加进来补偿，即 $\widehat{\Z}^{\times}\times\Q^{\times}\times\R^{\times}\to \mathbb{A}^{\times}_{\Q}$，则 $-1$ 在 $\Q^{\times}$ 与 $\R^{\times}$ 中被重复计入，所以必须把 $\R^{\times}$ 换成 $\R_{>0}$（见习题集 2 第 1 题）。

我们希望有一个阿代尔配对 $\mathbb{A}^{\times}_{\Q}\times\Q^{\times}\to\{1,-1\}$，从伊代尔与有理数映入 $\Z/2\Z$，它应当经过平方分解：
\[
\begin{CD}
\mathbb{A}^{\times}_{\Q}\times\Q^{\times} @>>> \{1,-1\}\\
@VVV @| \\
\Q^{\times}\backslash \mathbb{A}^{\times}_{\Q}/(\mathbb{A}^{\times}_{\Q})^2\times \Q^{\times}/(\Q^{\times})^2 @>>> \{1,-1\}
\end{CD}
\]
这里乘积左项中的那份 $\Q^{\times}$ 就是对角的「主伊代尔」，为了实现 (4.1)，这个配对也应当经过它分解。

这个映射应当诱导同构（某种意义上这正是「非退化性」）
\[
\Q^{\times}\backslash \mathbb{A}^{\times}_{\Q}/(\mathbb{A}^{\times}_{\Q})^2 \xrightarrow{\ \sim\ } \operatorname{Hom}\bigl(\Q^{\times}/(\Q^{\times})^2,\{1,-1\}\bigr),
\]
也就是说，这个映射不该恒等于 $1$，也不该是「什么古怪的东西」。

于是固定
\[
x=(x_p)_p\in \mathbb{A}^{\times}_{\Q},
\]
其中 $p$ 既可以是素数也可以是 $\infty$；并如下定义所求的配对：
\[
y\mapsto \prod_p (x_p,y)_p,
\]
这里我们把 $y$ 看作 $\Q_p$ 中的元素，而 $(\cdot,\cdot)_p$ 表示 $p$ 处的希尔伯特符号，即在 $\Q_p$ 上的希尔伯特符号。然而，这个无穷乘积是否良定义，先验地并不清楚；为此我们引入下面的引理：

\textbf{引理 4.3.} 对除有限多个 $p$ 之外的所有 $p$，$(x_p,y)_p=1$。

\textit{证明.} 事实上，对除有限多个 $p$ 之外的所有 $p$，都有 $p\neq 2,\infty$、$x_p\in\Z^{\times}_p$ 且 $y\in\Z^{\times}_p$；由驯符号的恒等式（因为 $x_p$ 与 $y$ 的赋值都是 $0$），这蕴含 $(x_p,y)_p=1$。\hfill$\square$

现在，这个映射当然是双乘性的（每一项都是），同样也当然地模平方分解，即它是映射
\[
\mathbb{A}^{\times}_{\Q}/(\mathbb{A}^{\times}_{\Q})^2\to\{1,-1\},
\]
因为若在任一侧乘上平方，希尔伯特符号不变。于是，这个映射在第一个因子上模 $\Q^{\times}$ 分解，当且仅当下面的断言成立：

\textbf{断言 4.4.} 对一切 $x,y\in\Q^{\times}$，都有 $\prod_p (x,y)_p=1$（也就是说，这个映射在第一个因子乘以 $\Q^{\times}$ 中的元素时不变；由双乘性，这意味着我们只是多提了一个这样的因子出来）。

\textbf{注 4.5.} 这对一切数域都成立（用一般的素位）。本讲中我们将看到，这（近似地）是二次互反律的一种重新包装。这一性质是素数之间的一种「共谋」：「$p$ 进域在幕后彼此通气；它们虽然彼此分离，却协力保证乘积为 $1$。」形容这类「共谋」的词就是「互反律」。

\textit{证明（断言 4.4）.} 首先，由于映射在乘以平方下不变，可以假设 $x=\pm p_1\cdots p_r$、$y=\pm q_1\cdots q_s$。于是双乘性蕴含我们可以取 $x\in\{-1,2,p\}$、$y\in\{-1,2,q\}$，其中 $p$ 与 $q$ 表示奇素数。我们只在 $p$ 与 $q$ 为互异的奇素数时证明该断言。

我们要证明
\[
(p,q)_{\infty}\times\prod_{\ell}(p,q)_{\ell}=1,
\]
其中 $\ell$ 遍历所有素数。第一项为 $1$，因为 $p$ 与 $q$ 都是正数；同样，$\ell$ 取异于 $p,q$ 的奇素数时也可以忽略。于是问题归结为证明
\[
(p,q)_2\cdot (p,q)_p\cdot (p,q)_q=1 .
\]

现在，正如习题集 1 第 2(b) 题所示，
\[
(p,q)_p=\left(\frac{(-1)^{v_p(p)v_p(q)}\,q^{v_p(p)}\big/p^{v_p(q)}}{p}\right)=\left(\frac{q}{p}\right),
\]
类似地 $(p,q)_q=\left(\frac{p}{q}\right)$。这里回忆
\[
\left(\frac{n}{p}\right):=\begin{cases} 1, & n \text{ 是模 } p \text{ 的平方},\\ -1, & \text{否则},\end{cases}
\]
其中 $n$ 与 $p$ 互素。此外，
\[
(p,q)_2:=(-1)^{\varepsilon(p)\varepsilon(q)}=(-1)^{\frac{p-1}{2}\cdot\frac{q-1}{2}},
\]
除非 $p\equiv q\equiv 3 \pmod 4$，它都等于 $1$。这样我们就归结到了初等的同余条件，而这正是二次互反律的陈述。\hfill$\square$

\textbf{注 4.6.} 二次互反律使得我们可以通过逐次归约来高效计算勒让德符号。

\textit{证明（二次互反律）.} 把勒让德符号看作映射
\[
\left(\frac{\cdot}{p}\right):\F^{\times}_p\to\{1,-1\},
\]
并把 $\F^{\times}_p$ 重新解释为 $\Gal(\Q(\zeta_p)/\Q)$，于是它是该 Galois 群唯一的非平凡二次特征标。这个特征标被 $\Q(\zeta_p)$（在 $\Q$ 上）唯一的二次子域所编码。我们希望把这个域写成 $\Q(\sqrt{d})$ 的形式。我们想要某个
\[
x=\sum_{n=1}^{p-1}x_n\zeta_p^{\,n}\in\Q(\zeta_p),
\]
使得对一切 $m\in(\Z/p\Z)^{\times}$，
\[
\sum_{n=1}^{p-1}x_n\zeta_p^{\,mn}=m\cdot x=\left(\frac{m}{p}\right)\cdot x=\sum_{n=1}^{p-1}\left(\frac{m}{p}\right)x_n\zeta_p^{\,n},
\]
这是作为 Galois 作用的等式：因为 $x_n\in\Q$ 不动，而作用为 $\zeta_p\mapsto\zeta_p^{\,m}$。也就是说，Galois 群把 $x$ 平移 $\pm 1$，这蕴含 $x\in\Q(\zeta_d)$，即我们的二次子域。
% [原文疑误] 此处原文印作 Q(ζ_d)；上文说二次子域形如 Q(√d)，Q(ζ_d) 疑为 Q(√d) 之误。照译并登记。

这个等式蕴含 $x_{mn}=\left(\frac{m}{p}\right)x_m$。反复解这个方程，我们最终得到元素（称为「高斯和」）
% [原文疑误] 原文此处印作 x_{mn}=(m/p)x_m；由下文 x_m=(m/p)x_1 及 Gauss 和的形式，右端下标应为 x_n。照译并登记。
\[
x=G:=\sum_{n=1}^{p-1}\left(\frac{n}{p}\right)\zeta_p^{\,n}.
\]
由构造我们知道 $G^2\in\Q$，但现在我们想知道它究竟是哪一个（事实上我们将看到，它是 $p$ 或 $-p$）。

设 $\chi:k^{\times}\to\C^{\times}$ 是乘法特征标，$\psi:k\to\C^{\times}$ 是加法特征标，其中 $k$ 是任意有限域。令
\[
G_{\chi,\psi}:=\sum_{x\in k^{\times}}\chi(x)\psi(x).
\]

\textbf{注 4.7.} 一个有趣的类比：伽马函数定义为
\[
\Gamma(\chi_s):=\int_{\R_{>0}}e^{-t}t^{s}\,\frac{dt}{t},
\]
这就像一个高斯和，其中取 $\psi(t):=e^{-t}$、$\chi_s(t):=t^{s}$。

我们需要下面的引理：

\textbf{引理 4.8.} 若 $\chi$ 与 $\psi$ 都不是平凡特征标，则 $G_{\chi,\psi}\cdot G_{\chi^{-1},\psi^{-1}}=\#k$。

\textit{证明.} 我们有
\[
\begin{aligned}
G_{\chi,\psi}\cdot G_{\chi^{-1},\psi^{-1}}
&=\sum_{x,y\in k^{\times}}\chi(x)\psi(x)\chi^{-1}(y)\psi^{-1}(y)\\
&=\sum_{x,y\in k^{\times}}\chi(x/y)\psi(x-y)\\
&=\sum_{z,y\in k^{\times}}\chi(z)\psi\bigl(y(z-1)\bigr),
\end{aligned}
\]
其中作了变量替换 $z:=x/y$，即 $x=zy$。现在，若 $z\neq 1$，则 $y(z-1)$ 取遍 $k^{\times}$ 中所有值，于是（由非退化性）$\sum_{w\in k}\psi(w)=0$ 成立。因此我们得到
\[
\begin{aligned}
G_{\chi,\psi}\cdot G_{\chi^{-1},\psi^{-1}}
&=\sum_{z\in k^{\times}}\chi(z)\sum_{y\in k^{\times}}\psi\bigl(y(z-1)\bigr)\\
&=\sum_{z\in k^{\times}\setminus\{1\}}\chi(z)\bigl(-\psi(0)\bigr)+\chi(1)\sum_{y\in k^{\times}}\psi(0)\\
&=\chi(1)+1\cdot\#k^{\times}=1+\#k^{\times}=\#k,
\end{aligned}
\]
因为类似地 $\sum_{w\in k^{\times}}\chi(w)=0$。\hfill$\square$

现在，当 $\psi$ 对应 $\zeta_p$ 的某个幂、$\chi$ 取乘法勒让德特征标时，我们想知道 $G_{\chi^{-1},\psi^{-1}}$ 是什么。我们有
\[
G_{\chi^{-1},\psi^{-1}}=\sum_{n=1}^{p-1}\left(\frac{n}{p}\right)\zeta_p^{-n}
=\sum_{n=1}^{p-1}\left(\frac{-n}{p}\right)\zeta_p^{\,n}
=\left(\frac{-1}{p}\right)\sum_{n=1}^{p-1}\left(\frac{n}{p}\right)\zeta_p^{\,n}
=\left(\frac{-1}{p}\right)G .
\]
于是
\[
G_{\chi,\psi}\cdot G_{\chi^{-1},\psi^{-1}}=G\cdot\left(\frac{-1}{p}\right)G=p,
\]
从而
\[
G^2=\left(\frac{-1}{p}\right)\cdot p=(-1)^{(p-1)/2}\cdot p,
\]
而 $G$ 是 $p$ 或 $-p$ 的平方根，视条件而定。

现在回忆：$\Q(\zeta_p)/\Q$ 在 $q\neq p$ 处非分歧，并且有同构
\[
\Gal(\Q(\zeta_p)/\Q)\cong(\Z/p\Z)^{\times},\qquad \mathrm{Frob}_q\mapsto q \bmod p .
\]
于是 $\left(\frac{q}{p}\right)=1$ 当且仅当 $\mathrm{Frob}_q$ 固定 $G$；事实上，
\[
\begin{aligned}
\mathrm{Frob}_q(G)
&=\mathrm{Frob}_q\Bigl(\sum_{n=1}^{p-1}\left(\frac{n}{p}\right)\zeta_p^{\,n}\Bigr)
=\sum_{n=1}^{p-1}\left(\frac{n}{p}\right)\zeta_p^{\,qn}
=\sum_{n=1}^{p-1}\left(\frac{q^{-1}n}{p}\right)\zeta_p^{\,n}\\
&=\sum_{n=1}^{p-1}\left(\frac{qn}{p}\right)\zeta_p^{\,n}
=\sum_{n=1}^{p-1}\left(\frac{q}{p}\right)\left(\frac{n}{p}\right)\zeta_p^{\,n}
=\left(\frac{q}{p}\right)G .
\end{aligned}
\]
% 原文此处的中间一步抽出为 "( n/q" 套在勒让德符号里，按 q 的逆元排作 (q^{-1}n/p)。

此外，
\[
G^{q-1}=(G^2)^{(q-1)/2}=\bigl((-1)^{(p-1)/2}\cdot p\bigr)^{(q-1)/2}\equiv (-1)^{\frac{p-1}{2}\cdot\frac{q-1}{2}}\left(\frac{p}{q}\right)\pmod q,
\]
所以
\[
\left(\frac{q}{p}\right)G^2=G^{q+1}\equiv (-1)^{\frac{p-1}{2}\cdot\frac{q-1}{2}}\left(\frac{p}{q}\right)G^2 \pmod q .
\]
除以 $G^2=(-1)^{(p-1)/2}\cdot p$（它模 $q$ 可逆）之后，得到
\[
\left(\frac{q}{p}\right)\equiv (-1)^{\frac{p-1}{2}\cdot\frac{q-1}{2}}\left(\frac{p}{q}\right)\pmod q,
\]
即
\[
\left(\frac{p}{q}\right)\left(\frac{q}{p}\right)=(-1)^{\frac{p-1}{2}\cdot\frac{q-1}{2}},
\]
如所愿。\hfill$\square$
